\documentclass[12pt,a4paper]{article}

\setlength{\topmargin}{0.0in}
\setlength{\oddsidemargin}{0.33in}
\setlength{\textheight}{9.0in}
\setlength{\textwidth}{6.0in}
\renewcommand{\baselinestretch}{1.25}



\usepackage{hyperref}
% \usepackage[utf8]{inputenc}
% \usepackage[T1]{fontenc}
\usepackage{float}
\usepackage{enumitem}
\usepackage{amsmath}
\usepackage{amssymb}
\usepackage{tikz}
\usetikzlibrary{decorations.pathreplacing}
\usepackage{amsfonts}
\usepackage{graphicx}
\usepackage{subcaption}
\usepackage{caption}
\newcommand{\dv}[2]{\frac{\mathrm{d} #1}{\mathrm{d} #2}}
\usepackage{caption}
\captionsetup[figure]{font=small} 
% \captionsetup[table]{font=small}
\usepackage[capitalize,nameinlink]{cleveref}
\usepackage{mathtools}
\newcommand{\defeq}{\vcentcolon=}
\newcommand{\eqdef}{=\vcentcolon}

% Define environments
\newtheorem{theorem}{Theorem}
\newtheorem{definition}{Definition}
\newtheorem{lemma}{Lemma}
\newtheorem{corollary}{Corollary}



\title{Logs from June $18^{\text{th}}$ to June $25^{\text{th}}$}
\author{}
\date{}
% \author{Maggie McCarthy}
% \date{May 2026}

\begin{document}

\maketitle

\section{Meeting June 18}


% \subsection{Summary of last week's work:}




% \begin{enumerate}
% \end{enumerate}



\subsection{Meeting Notes}

\begin{itemize}
    \item Could reach out to Matt to ask for a reference/proof for $\sigma_{\text{inf}}(B) = \sqrt{\sigma_{\text{inf}}(B^*B)}$ and for the Rayleigh quotient results in infinite dimensions. These results are both readily available in finite dimensions.
    \item Need the continuous basis to be a core - these is most likely a property of the Hilbert space (dont worry about it). 
    \item Do not worry to much about SCI. In general, we need three limits and in some cases we can get down to one.
    \item Focus on PseudoSpec instead of the minimization. For a general implementation later, we will force the user to input a $g$ sequence instead of having to choose it ourselves. 
    \item The spectral discretization works for the primal Maxwell problem (not mixed) because you have just one element so there is no continuity to enforce. 
    \item Next, we want to use DWR to improve $\Phi_n(z, A).$
    \item Next meeting is Thursday June 25 at 10 am.
\end{itemize}



\subsection{This Week's Plan}

\begin{enumerate}
    \item Work through some DWR literature in general and for Eigenproblems.
    \item Use DWR to compute an error estimare for the computed eigenvalue for the folded operator. This should allow us to attain a better $\Phi_n(z, A).$ Plot this to see!

\end{enumerate}


\section{Daily Logs}

\textbf{Thursday}

\begin{enumerate}
    \item Meeting with Patrick (notes above).
    \item Prepped for Boffi meeting.
    \item Boffi talk.
    \item More Boffi prep and meeting.
\end{enumerate}

\textbf{Friday}

\begin{enumerate}
    \item Started going through Ranacher text. Got through Ch.1 and most of Ch.2 My family arrived today for a visit!
\end{enumerate}

\textbf{Monday}

\begin{enumerate}
    \item Worked through bits of Chapter 2 and 3.
    \item Started Ch.4.
\end{enumerate}

\textbf{Tuesday}

\begin{enumerate}
    \item Worked through Ch.4 and Ch.6.
    \item Started Ch.7 (eigenvalue problems).
\end{enumerate}

\textbf{Wednesday}

\begin{enumerate}

\item Finished relevant parts of Ch.7 (7.1)
\item Attempted to `translate' the method of 7.1 to the Maxwell folded operator eigenvalue problem in Matt's algorithm (see notes).
\item Attempted an implementation - lots of questions and issues arose!
\end{enumerate}

\end{document}
